\section{The integrated certificate and its sign convention (main Proposition~\mref{prop:integrated})}\label{si:integrated}

\paragraph{Instance data.} Planar double integrator $\dot p=v$, $\dot v=u$, $u\in U=\conv\{\pm n_1,\pm n_2,\pm n_3\}$ with $n_1=(1,0)$, $n_2=(-\tfr35,\tfr45)$, $n_3=(-\tfr35,-\tfr45)$ (a hexagon with $h_U(\pm n_\ell)=1$); position facets $n_\ell^\top p\le2$; branch $\ell$ starts at $p_\ell(0)=\tfr{34}{25}n_\ell$, $v_\ell(0)=n_\ell$; the branch is revealed at $\tau$; horizon $T=\tau+1$. (The companion instance also carries velocity facets $|v_i|\le\tfr65$; they are not needed for the certificate direction and are omitted here.) The critical coordinate of branch $\ell$ is
\[
\zeta_\ell(T)=n_\ell^\top p_\ell(T)=\tfr{34}{25}+T+\int_0^T(T-\sigma)\,n_\ell^\top u_\ell(\sigma)\dd\sigma,
\]
so the terminal slack is $F_\ell=\beta_\ell-\int_0^Tk_\ell(\sigma)^\top u_\ell(\sigma)\dd\sigma$ with $\beta_\ell=\tfr{16}{25}-T$ and $k_\ell(\sigma)=(T-\sigma)n_\ell$.

\paragraph{The computation.} With $\lambda=(\tfr38,\tfr5{16},\tfr5{16})$: $\sum_\ell\lambda_\ell n_\ell=(\tfr38-\tfr3{16}-\tfr3{16},\ \tfr14-\tfr14)=(0,0)$, so the blind-window integrand $h_U(-(T-\sigma)\sum_\ell\lambda_\ell n_\ell)$ vanishes; after $\tau$ each branch is its own cell and $h_U(-(T-\sigma)n_\ell)=(T-\sigma)$; $\int_\tau^{T}(T-\sigma)\dd\sigma=\tfr12(T-\tau)^2=\tfr12$. Hence
\[
\Gamma_h(\tau)=\sum_\ell\lambda_\ell\beta_\ell+\sum_\ell\lambda_\ell\cdot\tfr12=\tfr{16}{25}-(\tau+1)+\tfr12=\tfr7{50}-\tau ,
\]
negative for $\tau>\tfr7{50}$, which is the obstruction (checks G1.1, G1.4). The script evaluates $\Gamma_h$ at $\tau\in\{0,\tfr1{10},\tfr7{50},\tfr15,\tfr12\}$ from the general formula (function \texttt{gamma\_h}, which integrates the kernels exactly and evaluates $h_U$ over the hexagon's vertices).

\paragraph{Sign convention.} Main \meqref{eq:gammah} evaluates $h_U$ at the \emph{negative} pooled kernel because $-g^\top u\le h_U(-g)$ is the bound used on each cell. On the hexagon $h_U(g)=h_U(-g)$ (central symmetry), and since $\sum_\ell\lambda_\ell n_\ell=0$ even an unbalanced evaluation would give the same number on this instance (check G1.2). The convention is not vacuous: for the translated set $U+(\tfr1{10},0)$ one has $h(-n_1)=\tfr9{10}$ and $h(n_1)=\tfr{11}{10}$, and with the unbalanced weights $(\tfr12,\tfr14,\tfr14)$ the two conventions give different values of $\Gamma_h$ (check G1.3). The companion paper's certificate $\Gamma(\tau)=\tau-\tfr7{50}$ is the value of a \emph{minimisation} programme (minimal maximal facet violation), positive when obstructed; the present $\Gamma_h$ is an upper bound on the worst terminal slack, negative when obstructed; $\Gamma_h=-\Gamma$ on the instance (check G1.4). Earlier versions of this manuscript wrote $h_U(+\,\cdot)$ in the integrated formula; the change is recorded in the change log.

\paragraph{Sharpness.} The policy ``hold $u=0$ on $[0,\tau)$, then $u=-n_\ell$ for one second'' keeps $\zeta_\ell\le\tfr{34}{25}+\tau+\tfr12$ on every branch, which is $\le2$ exactly when $\tau\le\tfr7{50}$ (companion paper, Section on the three-branch instance); so the certificate threshold is the viability threshold. The companion's velocity facets are respected by this policy ($|v|\le1<\tfr65$).
